\documentclass[11pt,a4paper]{article}
\usepackage[margin=25mm]{geometry}
\usepackage{amsmath,amssymb,amsthm,hyperref}
\setlength{\emergencystretch}{2em}
\newtheorem{lemma}{Lemma}
\newtheorem{theorem}[lemma]{Theorem}
\newcommand{\Sp}{\operatorname{Sp}}
\newcommand{\gen}[1]{\langle#1\rangle}
\title{A centralizer bound for symplectic strings in dimension four}
\author{Human Compression project}
\date{}
\begin{document}\maketitle
We prove an upper bound for connected strings of involutions in
characteristic two by determining two mutual centralizers. The proof
uses Jordan normal form and elementary matrix calculations.

\begin{lemma}\label{two}
Over an algebraically closed field $k$ of characteristic two, two
noncommuting involutions in $\mathrm{SL}_2(k)$ have no common invariant
line, and their common centralizer in $M_2(k)$ consists of scalars.
The centralizer in $\mathrm{SL}_2(k)$ of a single nonidentity involution
is abelian.
\end{lemma}
\begin{proof}
A nonidentity involution is unipotent with one fixed line. If two
involutions preserved the same line, both would be upper unipotent
in a suitable basis and would commute. Otherwise use their distinct
fixed lines as a basis. Their matrices are
\[
 u=\begin{pmatrix}1&\alpha\\0&1\end{pmatrix},\qquad
 v=\begin{pmatrix}1&0\\\beta&1\end{pmatrix},\qquad
 \alpha\beta\ne0.
\]
Commuting with $u$ forces a matrix to have the form
$\left(\begin{smallmatrix}c&d\\0&c\end{smallmatrix}\right)$;
commuting also with $v$ forces $d=0$.
In determinant one the first centralizer has $c^2=1$, hence $c=1$,
and is abelian.
\end{proof}

\begin{lemma}\label{centralizers}
Let $V$ be a four-dimensional symplectic space over $k$.
Suppose $A=\gen{a,b}$ and $D=\gen{e,f}$ are generated by
noncommuting involutions and commute elementwise with each other.
Then the full centralizers $C_{\Sp(V)}(A)$ and
$C_{\Sp(V)}(D)$ commute elementwise.
\end{lemma}
\begin{proof}
Put $h=ab$. Since $a,b$ do not commute, $h^2\ne1$.
The symplectic form pairs generalized eigenspaces only when their
eigenvalues are reciprocal. Thus the multiplicities of $\lambda$ and
$\lambda^{-1}$ agree, and the eigenvalue-one primary space is
nondegenerate. Also $aha=h^{-1}$.

First, if $h$ is not semisimple, its symplectic centralizer is abelian.
Here are all possibilities, including their justification in dimension four.
If all eigenvalues are one, a Jordan block has size at least three,
since otherwise $h^2=1$. The partition $3+1$ is impossible.
Indeed write $h x_1=x_1$, $h x_2=x_2+x_1$,
$h x_3=x_3+x_2$ and $h y=y$.
Preserving the pairings of $(x_1,x_3)$, $(x_2,x_3)$ and $(x_2,y)$
successively gives
$B(x_1,x_2)=B(x_1,x_3)=B(x_1,y)=0$.
Together with alternation this makes $x_1$ radical.
The only remaining case is one Jordan block of size four, whose
matrix centralizer consists of polynomials in that block and is abelian.

If a nonunit eigenvalue has a nontrivial Jordan block, reciprocal pairing
uses all four dimensions: there are two size-two blocks at
$\lambda,\lambda^{-1}$, $\lambda\ne1$.
A commuting symplectic map is determined by its action on the first
block; that action centralizes a size-two Jordan block and lies in
its abelian polynomial algebra. The other possibility has a size-two
unipotent block on the eigenvalue-one plane and simple eigenvalues
$\lambda,\lambda^{-1}$ on its perpendicular plane.
Its symplectic centralizer is a product of the abelian centralizer in
Lemma~\ref{two} and a one-dimensional torus. These exhaust the
nonsemisimple cases by reciprocal multiplicities.
Since $D$ is nonabelian and centralizes $h$, none can occur.

Thus $h$ is semisimple. Its centralizer can be nonabelian only for
one of the spectra
\[
 (\lambda,\lambda^{-1},1,1),\qquad
 (\lambda,\lambda,\lambda^{-1},\lambda^{-1}),
 \qquad \lambda\ne1.                                  \tag{1}
\]
Indeed four simple nonunit eigenvalues give only a torus, and four
unit eigenvalues would give $h=1$.

In the first case write $V=W\perp Z$, where $h|_Z=1$
and $W$ is its nonunit plane. The involution $a$ preserves these
planes and exchanges the two eigenlines in $W$.
A matrix commuting with both $a|_W,h|_W$ is scalar, and a symplectic
scalar in characteristic two is the identity.
Consequently
\[
 C_{\Sp(V)}(A)=1_W\times C_{\Sp(Z)}(a|_Z).
\]
If $a|_Z\ne1$, this group is abelian by Lemma~\ref{two}, impossible
because it contains $D$. Hence $a|_Z=b|_Z=1$ and this centralizer
is exactly $1_W\times\Sp(Z)$.
The two generators of $D$ are noncommuting involutions on $Z$.
Lemma~\ref{two} says that they have no common fixed vector and have
only scalar commuting matrices on $Z$. There are therefore no
intertwining blocks between their trivial action on $W$ and their
action on $Z$: the absence of a fixed covector follows as well from
the invariant symplectic form on $Z$.
It follows that $C_{\Sp(V)}(D)=\Sp(W)\times1_Z$.
The two asserted centralizers commute.

In the second case let $E,E'$ be the two eigenspaces of $h$.
Choose a basis of $E$ and its image under $a$ as basis of $E'$.
Then $a$ exchanges the two coordinate planes and the form matrix is
\[
                    \begin{pmatrix}0&S\\S&0\end{pmatrix},
\]
where $S$ is invertible and symmetric; symmetry follows from
invariance under $a$ and alternation in characteristic two.
The centralizer of $A$ consists exactly of
$\operatorname{diag}(T,T)$ with $T^{\mathsf T}ST=S$.
If $S$ is nonalternating, orthogonalization and square roots in $k$
put it into the form $I_2$. Its isometries are exactly
\[
             T=\begin{pmatrix}t&1+t\\1+t&t\end{pmatrix}
             \quad(t\in k),
\]
an abelian group. This is again impossible because it contains $D$.

Thus $S$ is alternating; normalize it to
$J=\left(\begin{smallmatrix}0&1\\1&0\end{smallmatrix}\right)$.
Then $C_{\Sp(V)}(A)=\{\operatorname{diag}(T,T):T\in\mathrm{SL}_2(k)\}$.
The two matrices representing $e,f$ on $E$ are noncommuting
involutions. By Lemma~\ref{two}, every block of an endomorphism
commuting with $D$ is scalar. Such an endomorphism is a two-by-two
block matrix $(c_{ij}I_2)$. It preserves the displayed form precisely
when $(c_{ij})^{\mathsf T}J(c_{ij})=J$.
Every such block matrix commutes with every $\operatorname{diag}(T,T)$.
This proves the lemma in the final case.
\end{proof}

\begin{theorem}\label{bound}
Over any field of characteristic two, a sequence of involutions in a
four-dimensional symplectic group satisfying
\[
 [\rho_i,\rho_j]=1\quad(|i-j|>1),\qquad
 [\rho_i,\rho_{i+1}]\ne1
\]
has length at most five.
\end{theorem}
\begin{proof}
Extend scalars to an algebraic closure, which preserves matrix equalities
and inequalities. If six generators existed, put
$A=\gen{\rho_0,\rho_1}$ and $D=\gen{\rho_4,\rho_5}$.
They satisfy Lemma~\ref{centralizers}.
But $\rho_2\in C(D)$ and $\rho_3\in C(A)$, so
$\rho_2$ and $\rho_3$ commute, a contradiction.
\end{proof}

\paragraph{Exact source interface and attribution.}
For $q=2^k\geq4$, every string C-marking of $\Sp(4,q)$ is connected:
a disconnected string would generate two nontrivial commuting normal
subgroups, impossible in this simple group. Thus Theorem~\ref{bound}
gives the upper bound five in Brooksbank--Ferrara--Leemans,
\emph{Orthogonal groups in characteristic 2 acting on polytopes of
high rank}, arXiv:1709.02219v2, Theorem 6.1.
Simplicity of $\Sp(4,q)$ for $q\geq4$ is a retained classical input
only for this final interface, not either lemma or Theorem~\ref{bound}.
The source's existence and small-field assertions are separate obligations.
This component proves no existence assertion or small-field classification.
\end{document}
